\section{Experimental Design}
\label{sec:experiments}

The protocol asks whether predicting additional joint recovery improves
complete-decision quality beyond an executed greedy response; whether
compatibility structure, occupancy supervision and warm information explain
any improvement; and how published solvers behave on public graph
optimization tasks. Published whole-graph
solvers and priorities within a shared recovery kernel answer complementary
questions. Performance tables pair each published method with its citation
and execution cost; implementation configurations are provided separately.

\paragraph{Controlled tasks and splits.}
Recovery menus cross 64/256 alternatives per menu, eight menus, three
topologies (Erd\H{o}s--R\'enyi, components, bipartite), densities
$0.08/0.20/0.45$, and cross-menu coupling $0.04$. Resource graphs cross
48/96 visibility passes with two/four/eight ground stations, eight
satellites and six alternatives per pass. Their conflicts represent
satellite and ground occupation. These 18 menu and six resource cells
define 72 training and 24 validation graphs, with graph-seed splits preceding
action/state extraction. All policy results use the complete 24-graph
development validation split. Previously inspected public inputs remain
exposed benchmark data.

\paragraph{Published whole-graph comparisons.}
The completed public comparison retains all 40 original unit-weight graphs:
16 DIMACS clique complements~\cite{johnson1996dimacs}, 18 SATLIB UF and six
CBS clause-conflict graphs~\cite{hoos2000satlib}, each with three common
initial schedules. The seven native configurations are CHILS-p1 and
CHILS-p16-c16~\cite{grossmann2025chils}, the two
$m^2$wis variants~\cite{grossmann2024mwis}, both Struction
configurations~\cite{gellner2021struction}, and
WeightedBR~\cite{lamm2019exact}. Their $0.5/2$\,s soft search workpoints
give 1,680 runs. The 16-thread CHILS configuration is reported separately
from single-thread configurations. Three released neural recipes add 360
runs: DIFUSCO-SAT with 50 steps and one/four
samples~\cite{sun2023difusco}, and COExpander-SAT-CM with one
sample (CM $1\times1$)~\cite{ma2025coexpander}. Raw and accepted cardinalities
are reported separately. Released-model
training overlap is unresolved, so this benchmark is a transfer comparison.

\paragraph{Numerical and application scope.}
The public benchmark uses the original unit rewards and conflict edges,
so released unit-MIS models are evaluated within their mathematical domain.
The controlled resource family retains satellite and station interval
conflicts and positive rewards. We do not extrapolate released unit-MIS
models to weighted contact scheduling. Large original contact sources and
weighted route benchmarks require separate numerical-domain and policy
transfer experiments; they are not part of the completed policy comparison.

\paragraph{Matched recovery controls.}
Every policy uses the same eleven commitment/release cells
$|C|\in\{0,2,4,8\}$, $|E|\in\{0,16,64\}$, excluding the no-op, and
fixed replacement cap of 256. The common paid prefix executes all three
weight/degree greedy sweeps and compares their feasible memberships with
$S\cap R$, yielding $W$. AnytimeGreedyPortfolio (greedy-only) returns this prefix;
Lower reuses it without repeating sweeps. Immediate, Upper, P1, Lower and
RoundRobin prioritize the same CHILS-p1 requests: one thread, seed 17,
10/50/200\,ms workpoints and at most eight calls. Upper uses a
clique-partition bound; P1 uses clique-load discounts. Unavailable cells remain
in coverage. All methods share
feasibility checks and incumbent preservation.

\paragraph{Learned controls and selection.}
Capacity and the FreeOccupancy, NoAux, NoWarm and CheapSummary controls
retain the same $W$, lower and upper bounds, and execution options.
NoWarm masks warm indicators and warm fraction while retaining scalar
$L=w(W)$. The label design clones two actual execution snapshots under
the revised warm semantics. Before fitting, the protocol requires measuring
native extra above the shared prefix across the full pool. Graph encoders
use width 32 and two layers; training fixes Adam $10^{-3}$,
four-graph batches and 40 epochs. All seeds 17/29/43 contribute;
each fit selects its first minimum validation
conditional regret. The strongest of all six classical controls at
$D=0.556$\,s is selected using the complete validation grid.

\paragraph{Quality, cost and statistics.}
Measured runs use Linux, an RTX 4080 and PyTorch 1.11/CUDA 11.3;
native solvers use one thread except the separately marked CHILS-p16-c16.
The total deadlines are 0.139, 0.556 and 2.221\,s. We normalize the
timely gain $g_\pi(D)$ by $\max\{1,w(S)\}$; a late return retains $S$.
Timing includes
proposals, scopes, features, bounds, inference, transfers, native execution,
validation/rescoring, synchronization and caller return. Common-prefix gain and
additional returned gain are reported separately from source availability.
Resident graph/model setup is separate; each decision starts with empty
caches. Native cold wall time includes export, startup, parsing and validation;
neural resident request time and once-only setup remain separate from
native soft stops and controller deadlines. They do not establish a
matched-budget speedup.

Fits and initial seeds are averaged within each graph before aggregation.
Controlled domains receive equal mass; public DIMACS/UF/CBS means receive
equal family weight. Failures, late returns and zero-opportunity graphs keep
their original mass; N/A remains distinct from executable failure.
Queries and repeated fits are not independent observations. Paired graph
bootstrap intervals describe development uncertainty, resampling graphs
within each domain and retaining equal domain weight. They do not constitute
an independent confirmation test.
